\documentclass{article}
\usepackage{amsmath, amssymb, amsthm}
\usepackage{geometry}
\geometry{a4paper, margin=1in}
\usepackage{hyperref}

\title{Linear Algebra Problem Sets: The Dimension of a Vector Space}
\author{}
\date{}

\begin{document}

\maketitle
\tableofcontents
\newpage

\section{Concept Checklist \& Question Mapping}
This section outlines all the core concepts, problem types, and tricks extracted from the provided material. Each concept is mapped to its corresponding questions in the problem set to ensure comprehensive coverage.

\begin{itemize}
    \item \textbf{Concept 1: Parametric Subspaces (Type 1)} 
    \begin{itemize}
        \item \textit{Description:} Factoring out parameters to express a generic vector as a linear combination of specific numerical vectors. Finding the basis and dimension by checking for linear independence.
        \item \textit{Mapped Questions:} 1, 2, 3, 4
    \end{itemize}
    
    \item \textbf{Concept 2: Subspace Spanned by Vectors (Type 2)}
    \begin{itemize}
        \item \textit{Description:} Finding the dimension of a span by placing vectors as columns in a matrix, row reducing, and counting the pivot columns. Utilizing the "Multiple" shortcut.
        \item \textit{Mapped Questions:} 5, 6, 7, 8
    \end{itemize}

    \item \textbf{Concept 3: Matrix Subspaces (Type 3)}
    \begin{itemize}
        \item \textit{Description:} Determining the dimensions of the Null Space (Nul A), Column Space (Col A), and Row Space (Row A) directly from a matrix using pivots and free variables.
        \item \textit{Mapped Questions:} 9, 10, 11, 12
    \end{itemize}

    \item \textbf{Concept 4: Rank Theorem Applications (Type 4)}
    \begin{itemize}
        \item \textit{Description:} Using the algebraic relation $\text{rank } A + \text{nullity } A = n$ (where $n$ is the number of columns). Understanding that $\text{rank } A = \text{rank } A^T$.
        \item \textit{Mapped Questions:} 13, 14, 15, 16
    \end{itemize}

    \item \textbf{Concept 5: Polynomial Spaces and Isomorphisms (Type 5)}
    \begin{itemize}
        \item \textit{Description:} Translating abstract vector spaces (like $\mathbb{P}_n$) into $\mathbb{R}^n$ using coordinate vectors relative to a standard basis. Using the Triangular Matrix Cheat.
        \item \textit{Mapped Questions:} 17, 18, 19, 20
    \end{itemize}

    \item \textbf{Concept 6: Conceptual True/False (Type 6)}
    \begin{itemize}
        \item \textit{Description:} Testing theoretical understanding of subspaces (must contain the origin), the Basis Theorem shortcut, and definitions of nullity and rank.
        \item \textit{Mapped Questions:} 21, 22, 23, 24
    \end{itemize}

    \item \textbf{Concept 7: Grassmann's Formula (Type 7)}
    \begin{itemize}
        \item \textit{Description:} Finding the dimension of the sum of two subspaces using $\dim(U + W) = \dim(U) + \dim(W) - \dim(U \cap W)$.
        \item \textit{Mapped Questions:} 25, 26, 27
    \end{itemize}

    \item \textbf{Concept 8: Find the Flaw (Type 8)}
    \begin{itemize}
        \item \textit{Description:} Diagnostic testing to identify common pitfalls: the RREF Column Trap, confusing rows and columns in the Rank Theorem, and forgetting the constant term in polynomials.
        \item \textit{Mapped Questions:} 28, 29, 30
    \end{itemize}
\end{itemize}

\newpage

\section{Problem Sets}

\subsection*{Type 1: Basis and Dimension of a Parametric Subspace}
\textbf{Problem 1:} Find a basis and state the dimension for the subspace of vectors of the form $\begin{bmatrix} 3s - t \\ s + 2t \\ -s \end{bmatrix}$ for $s,t \in \mathbb{R}$. \\

\textbf{Problem 2:} Find a basis and state the dimension for the subspace of vectors of the form $\begin{bmatrix} a + 2b - c \\ -a + b - 2c \\ 2a + 2c \\ 3b - 3c \end{bmatrix}$ for $a,b,c \in \mathbb{R}$. \\

\textbf{Problem 3:} Find a basis and state the dimension for the subspace of vectors $\begin{bmatrix} x \\ y \\ z \\ w \end{bmatrix}$ in $\mathbb{R}^4$ such that $x = a$, $y = b$, $z = c$, and $w = a+b+c$ for $a,b,c \in \mathbb{R}$. \\

\textbf{Problem 4:} Find a basis and state the dimension for the subspace of vectors of the form $\begin{bmatrix} a+b \\ b+c \\ c+d \\ d+a \end{bmatrix}$ for $a,b,c,d \in \mathbb{R}$. \\

\subsection*{Type 2: Dimension of a Subspace Spanned by Vectors}
\textbf{Problem 5:} Find the dimension of the subspace spanned by the vectors $\begin{bmatrix} 1 \\ 2 \\ 3 \end{bmatrix}, \begin{bmatrix} 2 \\ 4 \\ 6 \end{bmatrix}, \begin{bmatrix} 1 \\ 0 \\ 1 \end{bmatrix}$. \\

\textbf{Problem 6:} Find the dimension of the subspace spanned by the vectors $\begin{bmatrix} 1 \\ 1 \\ 1 \\ 1 \end{bmatrix}, \begin{bmatrix} 1 \\ 2 \\ 3 \\ 4 \end{bmatrix}, \begin{bmatrix} 2 \\ 3 \\ 4 \\ 5 \end{bmatrix}, \begin{bmatrix} 0 \\ 1 \\ 2 \\ 3 \end{bmatrix}$. \\

\textbf{Problem 7:} Find the dimension of the subspace spanned by the vectors $\begin{bmatrix} 1 \\ 0 \\ 0 \end{bmatrix}, \begin{bmatrix} 0 \\ 1 \\ 0 \end{bmatrix}, \begin{bmatrix} 0 \\ 0 \\ 1 \end{bmatrix}, \begin{bmatrix} 1 \\ 1 \\ 1 \end{bmatrix}, \begin{bmatrix} 2 \\ 3 \\ 4 \end{bmatrix}$. \\

\textbf{Problem 8:} Find the dimension of the subspace spanned by the vectors $\begin{bmatrix} 1 \\ -1 \\ 0 \\ 0 \end{bmatrix}, \begin{bmatrix} 0 \\ 0 \\ 1 \\ -1 \end{bmatrix}$. \\

\subsection*{Type 3: Dimensions of Matrix Subspaces}
\textbf{Problem 9:} Determine the dimensions of Nul $A$, Col $A$, and Row $A$ for the matrix $A = \begin{bmatrix} 1 & 2 & 0 & 4 \\ 0 & 0 & 1 & 5 \\ 0 & 0 & 0 & 0 \end{bmatrix}$. \\

\textbf{Problem 10:} Determine the dimensions of Nul $A$, Col $A$, and Row $A$ for the matrix $A = \begin{bmatrix} 1 & 1 & 1 \\ 2 & 2 & 2 \\ 3 & 3 & 3 \\ 4 & 4 & 4 \end{bmatrix}$. \\

\textbf{Problem 11:} Determine the dimensions of Nul $A$, Col $A$, and Row $A$ for the matrix $A = \begin{bmatrix} 1 & 0 & 0 \\ 0 & 2 & 0 \\ 0 & 0 & 3 \end{bmatrix}$. \\

\textbf{Problem 12:} Determine the dimensions of Nul $A$, Col $A$, and Row $A$ for the matrix $A = \begin{bmatrix} 1 & -2 & 3 & -4 & 5 \\ 0 & 1 & -1 & 2 & -2 \\ 0 & 0 & 0 & 1 & 1 \end{bmatrix}$. \\

\subsection*{Type 4: Applications of the Rank Theorem}
\textbf{Problem 13:} If a $6 \times 9$ matrix $A$ has rank 5, find nullity $A$, rank $A$, and rank $A^T$. \\

\textbf{Problem 14:} If the nullity of a $4 \times 7$ matrix $A$ is 3, what are the dimensions of the column and row spaces of $A$? \\

\textbf{Problem 15:} If an $8 \times 5$ matrix $A$ has a null space of dimension 2, find the rank of $A$. \\

\textbf{Problem 16:} If an $n \times n$ matrix $A$ has nullity 0, what is the rank of $A$? \\

\subsection*{Type 5: Polynomial Spaces and Isomorphisms}
\textbf{Problem 17:} What is the dimension of $\mathbb{P}_6$, the vector space of all polynomials of degree at most 6? \\

\textbf{Problem 18:} Show that the polynomials $1, 1+t, 1+t+t^2$ form a basis for $\mathbb{P}_2$. \\

\textbf{Problem 19:} Let $\mathcal{B} = \{1, t, t^2\}$ be the standard basis for $\mathbb{P}_2$. Find the coordinate vector of $p(t) = 4 + 3t - t^2$ relative to $\mathcal{B}$. \\

\textbf{Problem 20:} Determine if the polynomials $1+t, t+t^2, 1-t^2$ are linearly independent in $\mathbb{P}_2$. \\

\subsection*{Type 6: Conceptual True/False}
\textbf{Problem 21:} (T/F) A plane in $\mathbb{R}^4$ that does not pass through the origin is a two-dimensional subspace of $\mathbb{R}^4$. \\

\textbf{Problem 22:} (T/F) The nullity of a matrix $A$ is equal to the number of free variables in the homogeneous equation $Ax = 0$. \\

\textbf{Problem 23:} (T/F) If a set of 4 vectors spans a 4-dimensional vector space, it must be linearly independent. \\

\textbf{Problem 24:} (T/F) The dimension of the row space of $A$ is always equal to the dimension of the null space of $A$. \\

\subsection*{Type 7: Grassmann's Formula}
\textbf{Problem 25:} Let $U$ and $W$ be 4-dimensional subspaces of $\mathbb{R}^7$. What is the minimum possible dimension of their intersection, $U \cap W$? \\

\textbf{Problem 26:} Let $U$ and $W$ be 3-dimensional subspaces of $\mathbb{R}^5$. If they intersect in a 2-dimensional plane, what is the dimension of their sum, $U + W$? \\

\textbf{Problem 27:} Let $U$ and $W$ be 2-dimensional planes in $\mathbb{R}^4$. Is it mathematically possible for them to intersect only at the origin? Justify using Grassmann's Formula. \\

\subsection*{Type 8: Advanced Diagnostic Testing: "Find the Flaw"}
\textbf{Problem 28:} \textit{Find the flaw, explain why it is wrong, and provide the correct solution.} \\
\textbf{Flawed Solution:} "To find a basis for the column space of $A = \begin{bmatrix} 1 & 2 \\ 2 & 4 \end{bmatrix}$, we row reduce to the RREF matrix $\begin{bmatrix} 1 & 2 \\ 0 & 0 \end{bmatrix}$. The first column has a pivot, so the basis for Col $A$ is the first column of the RREF matrix: $\begin{bmatrix} 1 \\ 0 \end{bmatrix}$." \\

\textbf{Problem 29:} \textit{Find the flaw, explain why it is wrong, and provide the correct solution.} \\
\textbf{Flawed Solution:} "A $3 \times 5$ matrix $A$ has a rank of 2. By the Rank Theorem, $\text{rank } A + \text{nullity } A = \text{number of rows}$. So $2 + \text{nullity } A = 3$, meaning nullity $A = 1$." \\

\textbf{Problem 30:} \textit{Find the flaw, explain why it is wrong, and provide the correct solution.} \\
\textbf{Flawed Solution:} "The vector space $\mathbb{P}_3$ consists of polynomials up to degree 3. Since the highest power of $t$ is 3, the dimension of the space is 3."

\newpage

\section{Detailed Solutions}

\subsection*{Type 1: Basis and Dimension of a Parametric Subspace}
\textbf{Solution 1:} 
Decompose the generic vector into a linear combination of vectors with parameters $s$ and $t$ as weights:
$$ \begin{bmatrix} 3s - t \\ s + 2t \\ -s \end{bmatrix} = s \begin{bmatrix} 3 \\ 1 \\ -1 \end{bmatrix} + t \begin{bmatrix} -1 \\ 2 \\ 0 \end{bmatrix} $$
The subspace is spanned by $v_1 = \begin{bmatrix} 3 \\ 1 \\ -1 \end{bmatrix}$ and $v_2 = \begin{bmatrix} -1 \\ 2 \\ 0 \end{bmatrix}$. Since neither vector is a scalar multiple of the other, they are linearly independent. \\
\textbf{Final Answer:} The basis is $\left\{ \begin{bmatrix} 3 \\ 1 \\ -1 \end{bmatrix}, \begin{bmatrix} -1 \\ 2 \\ 0 \end{bmatrix} \right\}$. The dimension is 2. \\

\textbf{Solution 2:}
Decompose the vector into a linear combination of parameters $a, b, c$:
$$ a \begin{bmatrix} 1 \\ -1 \\ 2 \\ 0 \end{bmatrix} + b \begin{bmatrix} 2 \\ 1 \\ 0 \\ 3 \end{bmatrix} + c \begin{bmatrix} -1 \\ -2 \\ 2 \\ -3 \end{bmatrix} $$
Let these vectors be $v_1, v_2, v_3$. Notice that $v_1 - v_2 = \begin{bmatrix} 1-2 \\ -1-1 \\ 2-0 \\ 0-3 \end{bmatrix} = \begin{bmatrix} -1 \\ -2 \\ 2 \\ -3 \end{bmatrix} = v_3$. 
Because $v_3$ is a linear combination of $v_1$ and $v_2$, it is redundant and can be discarded. $v_1$ and $v_2$ are not multiples of each other, so they are linearly independent. \\
\textbf{Final Answer:} The basis is $\left\{ \begin{bmatrix} 1 \\ -1 \\ 2 \\ 0 \end{bmatrix}, \begin{bmatrix} 2 \\ 1 \\ 0 \\ 3 \end{bmatrix} \right\}$. The dimension is 2. \\

\textbf{Solution 3:}
Write the vector in parametric form:
$$ \begin{bmatrix} a \\ b \\ c \\ a+b+c \end{bmatrix} = a \begin{bmatrix} 1 \\ 0 \\ 0 \\ 1 \end{bmatrix} + b \begin{bmatrix} 0 \\ 1 \\ 0 \\ 1 \end{bmatrix} + c \begin{bmatrix} 0 \\ 0 \\ 1 \\ 1 \end{bmatrix} $$
These three vectors are clearly linearly independent (each has a 1 in a unique position among the first three rows). \\
\textbf{Final Answer:} The basis is $\left\{ \begin{bmatrix} 1 \\ 0 \\ 0 \\ 1 \end{bmatrix}, \begin{bmatrix} 0 \\ 1 \\ 0 \\ 1 \end{bmatrix}, \begin{bmatrix} 0 \\ 0 \\ 1 \\ 1 \end{bmatrix} \right\}$. The dimension is 3. \\

\textbf{Solution 4:}
Decompose the vector:
$$ a \begin{bmatrix} 1 \\ 0 \\ 0 \\ 1 \end{bmatrix} + b \begin{bmatrix} 1 \\ 1 \\ 0 \\ 0 \end{bmatrix} + c \begin{bmatrix} 0 \\ 1 \\ 1 \\ 0 \end{bmatrix} + d \begin{bmatrix} 0 \\ 0 \\ 1 \\ 1 \end{bmatrix} $$
Let these be $v_1, v_2, v_3, v_4$. Notice that $v_1 - v_2 + v_3 - v_4 = \begin{bmatrix} 1-1+0-0 \\ 0-1+1-0 \\ 0-0+1-1 \\ 1-0+0-1 \end{bmatrix} = \begin{bmatrix} 0 \\ 0 \\ 0 \\ 0 \end{bmatrix}$. 
Thus, the vectors are linearly dependent. Row reducing the matrix formed by these 4 vectors yields 3 pivot columns. \\
\textbf{Final Answer:} A basis is $\left\{ \begin{bmatrix} 1 \\ 0 \\ 0 \\ 1 \end{bmatrix}, \begin{bmatrix} 1 \\ 1 \\ 0 \\ 0 \end{bmatrix}, \begin{bmatrix} 0 \\ 1 \\ 1 \\ 0 \end{bmatrix} \right\}$. The dimension is 3. \\

\subsection*{Type 2: Dimension of a Subspace Spanned by Vectors}
\textbf{Solution 5:}
Let the vectors be $v_1, v_2, v_3$. Using the "Multiple" shortcut, observe that $v_2 = 2v_1$. Thus, $v_2$ is redundant. The remaining vectors $v_1$ and $v_3$ are not scalar multiples of each other, so they are linearly independent. \\
\textbf{Final Answer:} The dimension is 2. \\

\textbf{Solution 6:}
Form a matrix $A$ with these vectors as columns and row reduce:
$$ A = \begin{bmatrix} 1 & 1 & 2 & 0 \\ 1 & 2 & 3 & 1 \\ 1 & 3 & 4 & 2 \\ 1 & 4 & 5 & 3 \end{bmatrix} \sim \begin{bmatrix} 1 & 1 & 2 & 0 \\ 0 & 1 & 1 & 1 \\ 0 & 2 & 2 & 2 \\ 0 & 3 & 3 & 3 \end{bmatrix} \sim \begin{bmatrix} 1 & 1 & 2 & 0 \\ 0 & 1 & 1 & 1 \\ 0 & 0 & 0 & 0 \\ 0 & 0 & 0 & 0 \end{bmatrix} $$
There are 2 pivot columns. \\
\textbf{Final Answer:} The dimension is 2. \\

\textbf{Solution 7:}
The subspace is spanned by 5 vectors in $\mathbb{R}^3$. The maximum possible dimension for any subspace in $\mathbb{R}^3$ is 3. Since the first three vectors are the standard basis vectors $e_1, e_2, e_3$, they alone span all of $\mathbb{R}^3$. \\
\textbf{Final Answer:} The dimension is 3. \\

\textbf{Solution 8:}
There are 2 vectors. Since neither is a scalar multiple of the other (one has zeros where the other has non-zeros), they are linearly independent. \\
\textbf{Final Answer:} The dimension is 2. \\

\subsection*{Type 3: Dimensions of Matrix Subspaces}
\textbf{Solution 9:}
The matrix is already in echelon form. There are pivot positions in row 1, col 1 and row 2, col 3.
\begin{itemize}
    \item Number of pivot columns = 2. Thus, $\dim(\text{Col } A) = 2$.
    \item By the Rank Theorem, $\dim(\text{Row } A) = \dim(\text{Col } A) = 2$.
    \item The matrix has $n = 4$ columns. $\dim(\text{Nul } A) = n - \text{rank} = 4 - 2 = 2$.
\end{itemize}
\textbf{Final Answer:} $\dim(\text{Nul } A) = 2$, $\dim(\text{Col } A) = 2$, $\dim(\text{Row } A) = 2$. \\

\textbf{Solution 10:}
Notice that rows 2, 3, and 4 are all scalar multiples of row 1. Row reducing will leave only 1 non-zero row.
\begin{itemize}
    \item Number of pivot columns = 1. Thus, $\dim(\text{Col } A) = 1$.
    \item $\dim(\text{Row } A) = 1$.
    \item The matrix has $n = 3$ columns. $\dim(\text{Nul } A) = 3 - 1 = 2$.
\end{itemize}
\textbf{Final Answer:} $\dim(\text{Nul } A) = 2$, $\dim(\text{Col } A) = 1$, $\dim(\text{Row } A) = 1$. \\

\textbf{Solution 11:}
The matrix is a diagonal matrix with non-zero entries on the main diagonal. It is already in echelon form with 3 pivots.
\begin{itemize}
    \item $\dim(\text{Col } A) = 3$.
    \item $\dim(\text{Row } A) = 3$.
    \item $n = 3$ columns. $\dim(\text{Nul } A) = 3 - 3 = 0$.
\end{itemize}
\textbf{Final Answer:} $\dim(\text{Nul } A) = 0$, $\dim(\text{Col } A) = 3$, $\dim(\text{Row } A) = 3$. \\

\textbf{Solution 12:}
The matrix is in echelon form. The pivots are in columns 1, 2, and 4.
\begin{itemize}
    \item Number of pivot columns = 3. Thus, $\dim(\text{Col } A) = 3$.
    \item $\dim(\text{Row } A) = 3$.
    \item $n = 5$ columns. $\dim(\text{Nul } A) = 5 - 3 = 2$.
\end{itemize}
\textbf{Final Answer:} $\dim(\text{Nul } A) = 2$, $\dim(\text{Col } A) = 3$, $\dim(\text{Row } A) = 3$. \\

\subsection*{Type 4: Applications of the Rank Theorem}
\textbf{Solution 13:}
\begin{itemize}
    \item We are given $\text{rank } A = 5$.
    \item The matrix has $n = 9$ columns. By the Rank Theorem, $\text{rank } A + \text{nullity } A = n \implies 5 + \text{nullity } A = 9 \implies \text{nullity } A = 4$.
    \item The rank of $A^T$ is the dimension of the row space of $A$, which equals the rank of $A$. Thus, $\text{rank } A^T = 5$.
\end{itemize}
\textbf{Final Answer:} nullity $A = 4$, rank $A = 5$, rank $A^T = 5$. \\

\textbf{Solution 14:}
The matrix has $n = 7$ columns. Using the Rank Theorem: $\dim(\text{Col } A) + \text{nullity } A = n$.
$$ \dim(\text{Col } A) + 3 = 7 \implies \dim(\text{Col } A) = 4 $$
Since the dimension of the row space always equals the dimension of the column space: $\dim(\text{Row } A) = 4$. \\
\textbf{Final Answer:} The dimension of the column space is 4, and the dimension of the row space is 4. \\

\textbf{Solution 15:}
The matrix has $n = 5$ columns. We are given $\dim(\text{Nul } A) = 2$.
$$ \text{rank } A + 2 = 5 \implies \text{rank } A = 3 $$
\textbf{Final Answer:} The rank of $A$ is 3. \\

\textbf{Solution 16:}
The matrix has $n$ columns. We are given nullity $A = 0$.
$$ \text{rank } A + 0 = n \implies \text{rank } A = n $$
\textbf{Final Answer:} The rank of $A$ is $n$ (the matrix is full rank). \\

\subsection*{Type 5: Polynomial Spaces and Isomorphisms}
\textbf{Solution 17:}
The standard basis for $\mathbb{P}_6$ is $\{1, t, t^2, t^3, t^4, t^5, t^6\}$. Counting these basis elements, there are 7. \\
\textbf{Final Answer:} The dimension of $\mathbb{P}_6$ is 7. \\

\textbf{Solution 18:}
Map each polynomial to its coordinate vector in $\mathbb{R}^3$ relative to the standard basis $\mathcal{B} = \{1, t, t^2\}$:
$$ p_1 = 1 \implies [p_1]_\mathcal{B} = \begin{bmatrix} 1 \\ 0 \\ 0 \end{bmatrix}, \quad p_2 = 1+t \implies [p_2]_\mathcal{B} = \begin{bmatrix} 1 \\ 1 \\ 0 \end{bmatrix}, \quad p_3 = 1+t+t^2 \implies [p_3]_\mathcal{B} = \begin{bmatrix} 1 \\ 1 \\ 1 \end{bmatrix} $$
Form a matrix with these vectors as columns:
$$ M = \begin{bmatrix} 1 & 1 & 1 \\ 0 & 1 & 1 \\ 0 & 0 & 1 \end{bmatrix} $$
Using the Triangular Matrix Cheat, this matrix is upper triangular with no zeros on the main diagonal. Therefore, it has 3 pivot positions, and the columns are linearly independent. By the Basis Theorem, 3 linearly independent vectors in a 3-dimensional space ($\mathbb{P}_2$) automatically form a basis. \\

\textbf{Solution 19:}
The polynomial is $p(t) = 4(1) + 3(t) - 1(t^2)$. The coefficients directly form the coordinate vector. \\
\textbf{Final Answer:} The coordinate vector is $[p]_\mathcal{B} = \begin{bmatrix} 4 \\ 3 \\ -1 \end{bmatrix}$. \\

\textbf{Solution 20:}
Map to coordinate vectors:
$$ [p_1]_\mathcal{B} = \begin{bmatrix} 1 \\ 1 \\ 0 \end{bmatrix}, \quad [p_2]_\mathcal{B} = \begin{bmatrix} 0 \\ 1 \\ 1 \end{bmatrix}, \quad [p_3]_\mathcal{B} = \begin{bmatrix} 1 \\ 0 \\ -1 \end{bmatrix} $$
Notice that $[p_1]_\mathcal{B} - [p_2]_\mathcal{B} = \begin{bmatrix} 1 \\ 1 \\ 0 \end{bmatrix} - \begin{bmatrix} 0 \\ 1 \\ 1 \end{bmatrix} = \begin{bmatrix} 1 \\ 0 \\ -1 \end{bmatrix} = [p_3]_\mathcal{B}$.
Since one vector is a linear combination of the others, they are linearly dependent. \\
\textbf{Final Answer:} They are linearly dependent. \\

\subsection*{Type 6: Conceptual True/False}
\textbf{Solution 21:} \textbf{False.} A plane is a subspace if and only if it passes through the origin (contains the zero vector). If it does not, it is not a vector space at all. \\

\textbf{Solution 22:} \textbf{True.} The columns of $A$ that are not pivot columns correspond exactly to the free variables in $Ax = 0$. The number of free variables dictates the number of basis vectors for the null space, which is the definition of nullity. \\

\textbf{Solution 23:} \textbf{True.} This is the Basis Theorem Shortcut. If you know the dimension of a space is $p$ (here, 4), and you are given exactly $p$ vectors that span the space, they are automatically linearly independent and form a basis. \\

\textbf{Solution 24:} \textbf{False.} The dimension of the row space equals the rank of $A$. The dimension of the null space equals $n - \text{rank } A$. They are only equal if $\text{rank } A = n/2$, which is not always true. \\

\subsection*{Type 7: Grassmann's Formula}
\textbf{Solution 25:}
By Grassmann's Formula: $\dim(U + W) = \dim(U) + \dim(W) - \dim(U \cap W)$.
Since $U$ and $W$ are subspaces of $\mathbb{R}^7$, their sum $U + W$ must also be a subspace of $\mathbb{R}^7$. Therefore, the maximum possible dimension for $U + W$ is 7.
$$ \dim(U + W) \le 7 $$
$$ 4 + 4 - \dim(U \cap W) \le 7 $$
$$ 8 - \dim(U \cap W) \le 7 \implies \dim(U \cap W) \ge 1 $$
\textbf{Final Answer:} The minimum possible dimension of their intersection is 1. \\

\textbf{Solution 26:}
We are given $\dim(U) = 3$, $\dim(W) = 3$, and $\dim(U \cap W) = 2$.
Using Grassmann's Formula:
$$ \dim(U + W) = 3 + 3 - 2 = 4 $$
\textbf{Final Answer:} The dimension of their sum is 4. \\

\textbf{Solution 27:}
If they intersect only at the origin, then $\dim(U \cap W) = 0$.
Using Grassmann's Formula:
$$ \dim(U + W) = 2 + 2 - 0 = 4 $$
Since $U + W$ is a subspace of $\mathbb{R}^4$, its maximum dimension is 4. Because $4 \le 4$ is a valid statement, it is mathematically possible for two 2D planes in $\mathbb{R}^4$ to intersect only at the origin. \\
\textbf{Final Answer:} Yes, it is possible. \\

\subsection*{Type 8: Advanced Diagnostic Testing: "Find the Flaw"}
\textbf{Solution 28:}
\textbf{Error:} The student fell into the "RREF Column Trap." They used the column from the row-reduced matrix instead of the original matrix. Row operations change the column space! \\
\textbf{Correct Solution:} The pivot is indeed in the first column. However, we must take the first column of the \textit{original} matrix $A$. The correct basis is $\left\{ \begin{bmatrix} 1 \\ 2 \end{bmatrix} \right\}$. \\

\textbf{Solution 29:}
\textbf{Error:} The student confused rows and columns in the Rank Theorem. The theorem uses $n$ (the number of columns), not $m$ (the number of rows). Nullity is about the domain ($\mathbb{R}^n$). \\
\textbf{Correct Solution:} The Rank Theorem states $\text{rank } A + \text{nullity } A = n$. The matrix has $n = 5$ columns. Thus, $2 + \text{nullity } A = 5 \implies \text{nullity } A = 3$. \\

\textbf{Solution 30:}
\textbf{Error:} The student forgot that polynomials include a constant term (degree 0). \\
\textbf{Correct Solution:} The standard basis for $\mathbb{P}_3$ is $\{1, t, t^2, t^3\}$. Counting these basis elements, there are exactly 4. The correct dimension is 4.

\newpage

\section{Concept Checklist \& Question Mapping (Reference)}
\begin{itemize}
    \item \textbf{Concept 1: Parametric Subspaces (Type 1)} - Questions 1, 2, 3, 4
    \item \textbf{Concept 2: Subspace Spanned by Vectors (Type 2)} - Questions 5, 6, 7, 8
    \item \textbf{Concept 3: Matrix Subspaces (Type 3)} - Questions 9, 10, 11, 12
    \item \textbf{Concept 4: Rank Theorem Applications (Type 4)} - Questions 13, 14, 15, 16
    \item \textbf{Concept 5: Polynomial Spaces and Isomorphisms (Type 5)} - Questions 17, 18, 19, 20
    \item \textbf{Concept 6: Conceptual True/False (Type 6)} - Questions 21, 22, 23, 24
    \item \textbf{Concept 7: Grassmann's Formula (Type 7)} - Questions 25, 26, 27
    \item \textbf{Concept 8: Find the Flaw (Type 8)} - Questions 28, 29, 30
\end{itemize}

\end{document}